\section*{Bibliographie}\label{I:bib}
\begin{description}
\item[{[1]}] M. ARTIN. Grothendieck topologies (notes miméographiées), Harvard University, 1962.
\item[{[2]}] M. ARTIN. Algebraic approximation of structures over complete local rings, Pub. Math. I.H.E.S., n° 36, \var{1959}{1969}{S45-ASTRA-A-0029}, p. \var{25--58}{23--58}{S45-ASTRA-A-0030}.
\item[{[3]}] M. ARTIN. Théorèmes de représentabilité pour les espaces algébriques (chap. VII, le théorème de finitude en cohomologie étale), Presses de l'Université de Montréal, 1973.
\item[{[4]}] M. ARTIN, A. GROTHENDIECK, J.L. VERDIER. Théorie des Topos et Cohomologie Étale des Schémas, SGA 4, Lecture notes in Maths., n° 269, 270 et 305, Springer, 1972/73.
\item[{[5]}] J. GIRAUD. Analysis situs (d'après Artin et Grothendieck). Séminaire Bourbaki, 1962/63, exp. \var{266}{256}{S45-ASTRA-A-0033}.
\item[{[6]}] A. GROTHENDIECK. Technique de descente et théorèmes d'existence en géométrie algébrique. I. Généralités. Descente par morphismes fidèlement plats. Séminaire Bourbaki, 1959/60, exp. 190.
\item[{[7]}] A. GROTHENDIECK. Revêtements Étales et Groupe Fondamental, SGA 1, Lecture notes in Maths., n° 224, Springer, 1971.
\item[{[8]}] A. GROTHENDIECK et J. DIEUDONNÉ. Éléments de géométrie algébrique, Pub. Math. I.H.E.S., n° 4, 8, 11, 17, 20, 24, 28, 32, 1960/67 [cité EGA].
\item[{[9]}] M. NAGATA. Imbedding of an abstract variety in a complete variety, J. Math. Kyoto Univ., 2-1, 1962, p. 1--10.
\item[{[10]}] M. NAGATA. Generalisation of the imbedding problem, J. Math. Kyoto Univ., 3-1, 1963, p. 89--102.
\item[{[11]}] M. RAYNAUD. Anneaux Locaux Henséliens, Lecture notes in Maths., n° 169, Springer, 1970.
\item[{[12]}] J.-P. SERRE. Géométrie Algébrique et Géométrie Analytique, Ann. de l'Institut Fourier, t. VI, 1956, p. \var{1--12}{1--42}{S45-ASTRA-A-0031} (cité GAGA).
\item[{[13]}] J.-P. SERRE. Corps Locaux, Paris, Hermann, 1968.
\item[{[14]}] J.-P. SERRE. Cohomologie galoisienne, Lecture notes in Maths., n° 5, Springer, 1965.
\item[{[15]}] J.-P. SERRE. Groupes algébriques et corps de classes, Paris, Hermann, 1959.
\item[{[16]}] C. TSEN. Divisionsalgebren über Funktionenkörper, Nach. Ges. Wiss. Göttingen, 1933, p. 335--339.
\end{description}
